% Chapter 4, Section 3 _Linear Algebra_ Jim Hefferon
%  http://joshua.smcvt.edu/linearalgebra
%  2001-Jun-12
\section{Rumus Laplace}
\textit{Bagian ini bersifat opsional.
  Satu-satunya bagian selanjutnya yang bergantung pada materi ini
  adalah Lima.III.}

Determinan merupakan sumber berbagai rumus yang menarik dan menghibur.
Berikut salah satu rumus yang sering digunakan
untuk menghitung determinan dengan tangan.



\subsectionoptional{Ekspansi Laplace}
\index{ekspansi determinan Laplace|(}

%<*dis:LaplaceIntro0>
Dalam ekspansi permutasi $\deter{T}$,
setiap suku memiliki tepat satu entri dari setiap baris dan kolom~$T$.
\begin{equation*}
  \deter{T}=\!\!
  \sum_{\text{permutasi\ }\phi}\!\!\!\!
     t_{1,\phi(1)}t_{2,\phi(2)}\cdots t_{n,\phi(n)}
                                 \deter{P_{\phi}}
\end{equation*}
Berikut bentuk rumus tersebut untuk matriks $\nbyn{3}$.
\begin{align*}
  \begin{vmat}
    t_{1,1}  &t_{1,2}  &t_{1,3} \\
    t_{2,1}  &t_{2,2}  &t_{2,3} \\
    t_{3,1}  &t_{3,2}  &t_{3,3} 
  \end{vmat}
  &=
  \begin{aligned}[t]
    &t_{1,1}t_{2,2}t_{3,3}\deter{P_{\phi_1}}
     +t_{1,1}t_{2,3}t_{3,2}\deter{P_{\phi_2}}
     +t_{1,2}t_{2,1}t_{3,3}\deter{P_{\phi_3}} \\
    &\hspace*{0em}
     +t_{1,2}t_{2,3}t_{3,1}\deter{P_{\phi_4}}
     +t_{1,3}t_{2,1}t_{3,2}\deter{P_{\phi_5}}
     +t_{1,3}t_{2,2}t_{3,1}\deter{P_{\phi_6}}
  \end{aligned}                                      \\
  &=
  \begin{aligned}[t]
    &t_{1,1}t_{2,2}t_{3,3}
     -t_{1,1}t_{2,3}t_{3,2}
     -t_{1,2}t_{2,1}t_{3,3}  \\
    &\hspace*{0em}
     +t_{1,2}t_{2,3}t_{3,1} 
     +t_{1,3}t_{2,1}t_{3,2}
     -t_{1,3}t_{2,2}t_{3,1}
  \end{aligned}
\end{align*}
%</dis:LaplaceIntro0>
%<*dis:LaplaceIntro1>
Kita dapat berfokus pada sebarang baris atau kolom; di sini kita
mengeluarkan entri-entri dari baris pertama $T$.
\begin{align*}
  &=
    t_{1,1}\cdot\big[t_{2,2}t_{3,3}
                      -t_{2,3}t_{3,2}\big]  \\
    &\quad
     -t_{1,2}\cdot\big[t_{2,1}t_{3,3}
                       -t_{2,3}t_{3,1}\big]  \\
    &\quad
     +t_{1,3}\cdot\big[t_{2,1}t_{3,2}
                       -t_{2,2}t_{3,1}\big]
  \tag{$*$}
\end{align*}
%</dis:LaplaceIntro1>
%<*dis:LaplaceIntro2>
Perhatikan bahwa di dalam pasangan kurung siku pertama terdapat
determinan
\begin{equation*}
  \begin{vmat}
    t_{2,2}  &t_{2,3} \\
    t_{3,2}  &t_{3,3}
  \end{vmat}
  \tag{$**$}
\end{equation*}
dan demikian pula, di dalam pasangan kurung lainnya terdapat determinan lain.
\begin{equation*}
  \begin{vmat}
    t_{2,1}  &t_{2,3} \\
    t_{3,1}  &t_{3,3}
  \end{vmat}
  \qquad
  \begin{vmat}
    t_{2,1}  &t_{2,2} \\
    t_{3,1}  &t_{3,2}
  \end{vmat}
  \tag{$*\!*\!*$}
\end{equation*}
%</dis:LaplaceIntro2>
%<*dis:LaplaceIntro3>
Penjelasan bagi determinan-determinan ini terletak pada permutasinya.
Jika kita menetapkan suatu baris dan kolom\Dash misalnya ketika meninjau semua
suku yang memuat~$t_{1,1}$, semua suku yang memuat~$t_{1,2}$, atau semua suku
yang memuat~$t_{1,3}$\Dash maka yang tersisa bagi permutasi adalah memilih satu
entri dari setiap baris dan kolom lainnya.

Matriks-matriks di bawah ini memberikan ilustrasi.
Di sebelah kiri, baris dan kolom yang memuat~$t_{1,1}$ diberi bayangan, dan
yang tersisa adalah matriks dari~($**$) di atas.
Demikian pula, pada dua matriks lainnya, yang tersisa adalah matriks-matriks
dari~($*\!*\!*$).
\definecolor{shadematrix}{gray}{0.8}
\begin{equation*}
  \left(\begin{array}{>{\columncolor{shadematrix}}ccc}
    \rowcolor{shadematrix}
    t_{1,1}  &t_{1,2}  &t_{1,3} \\
    t_{2,1}  &t_{2,2}  &t_{2,3} \\
    t_{3,1}  &t_{3,2}  &t_{3,3} 
  \end{array}\right)
  \quad
  \left(\begin{array}{c>{\columncolor{shadematrix}}cc}
    \rowcolor{shadematrix}
    t_{1,1}  &t_{1,2}  &t_{1,3} \\
    t_{2,1}  &t_{2,2}  &t_{2,3} \\
    t_{3,1}  &t_{3,2}  &t_{3,3} 
  \end{array}\right)
  \quad
  \left(\begin{array}{cc>{\columncolor{shadematrix}}c}
    \rowcolor{shadematrix}
    t_{1,1}  &t_{1,2}  &t_{1,3} \\
    t_{2,1}  &t_{2,2}  &t_{2,3} \\
    t_{3,1}  &t_{3,2}  &t_{3,3} 
  \end{array}\right)
\end{equation*}
%</dis:LaplaceIntro3>

%<*dis:LaplaceIntro4>
Dengan demikian, kita telah menjelaskan kemunculan determinan matriks-matriks
berukuran lebih kecil dalam persamaan~($*$).
Untuk menjelaskan tanda minus pada baris kedua persamaan tersebut,
tinjau bentuk berikut dari
$\deter{T}=t_{1,1}t_{2,2}t_{3,3}\deter{P_{\phi_1}}+\cdots+t_{1,3}t_{2,2}t_{3,1}\deter{P_{\phi_6}}$.
\begin{equation*}
  \begin{aligned}[t]
      &t_{1,1}\cdot \bigg[t_{2,2}t_{3,3}\begin{vmat}[r]
                                 1  &0  &0  \\
                                 0  &1  &0  \\
                                 0  &0  &1
                                \end{vmat}
                 +t_{2,3}t_{3,2}\begin{vmat}[r]
                                  1  &0  &0  \\
                                  0  &0  &1  \\
                                  0  &1  &0
                                \end{vmat}\,\bigg]    \\
         &\hbox{}\quad\hbox{}
          +t_{1,2}\cdot \bigg[t_{2,1}t_{3,3}\begin{vmat}[r]
                                        0  &1  &0  \\
                                        1  &0  &0  \\
                                        0  &0  &1
                                       \end{vmat}
                        +t_{2,3}t_{3,1}\begin{vmat}[r]
                                         0  &1  &0  \\
                                         0  &0  &1  \\
                                         1  &0  &0
                                        \end{vmat}\,\bigg]  \\
        &\hbox{}\quad\hbox{}
         +t_{1,3}\cdot \bigg[t_{2,1}t_{3,2}\begin{vmat}[r]
                                        0  &0  &1  \\
                                        1  &0  &0  \\
                                        0  &1  &0
                                       \end{vmat}
                        +t_{2,2}t_{3,1}\begin{vmat}[r]
                                         0  &0  &1  \\
                                         0  &1  &0  \\
                                         1  &0  &0
                                        \end{vmat}\,\bigg]
    \end{aligned}
\end{equation*}
Matriks $P_{\phi_1}$ dan~$P_{\phi_2}$ pada baris pertama sudah tersusun untuk
determinan~($**$), karena baris pertamanya sama seperti pada matriks identitas,
sedangkan dua baris lainnya tampak seperti baris-baris determinan tersebut.
Matriks $P_{\phi_3}$ dan~$P_{\phi_4}$ pada baris kedua dapat disusun dengan
cara yang sama melalui satu pertukaran baris, yang menempatkan
$\rowvec{1 &0 &0}$ pada baris pertama.
Tentu saja, pertukaran baris tersebut memberi tanda negatif pada baris itu.
Terakhir, menyusun matriks-matriks pada baris ketiga dengan cara yang sama
memerlukan dua pertukaran, sehingga tandanya tetap positif.
%</dis:LaplaceIntro4>

Singkatnya, kita memperoleh rumus berikut.
\begin{equation*}
  \deter{T}
  =t_{1,1}\cdot \begin{vmat}
            t_{2,2}  &t_{2,3}  \\
            t_{3,2}  &t_{3,3}
          \end{vmat}
   -t_{1,2}\cdot \begin{vmat}
             t_{2,1}  &t_{2,3}  \\
             t_{3,1}  &t_{3,3}
           \end{vmat}
   +t_{1,3}\cdot \begin{vmat}
             t_{2,1}  &t_{2,2}  \\
             t_{3,1}  &t_{3,2}
           \end{vmat}
\end{equation*}
Rumus dalam \nearbytheorem{th:LaPlaceExp}, yang menggeneralisasi hasil ini,
merupakan suatu \definend{rekurensi}\index{rekurensi}, karena determinan
dinyatakan sebagai kombinasi determinan-determinan.
Rumus ini tidak bersifat melingkar karena menyatakan determinan $\nbyn{n}$
dalam bentuk matriks-matriks berukuran lebih kecil.

\begin{definition} \label{df:Minor}
%<*df:Minor>
Untuk sebarang matriks \( \nbyn{n} \), \( T \), matriks \( \nbyn{(n-1)} \)
yang diperoleh dengan menghapus baris~\( i \) dan kolom~\( j \) dari \( T \)
adalah \definend{minor} \( i,j \)\index{minor, matriks}%
\index{determinan!minor}
\index{matriks!minor}
dari \( T \).
\definend{Kofaktor} \( i,j \)\index{kofaktor}
\index{determinan!menggunakan kofaktor}
% \index{matrix!cofactor} 
dari \( T \), yaitu \( T_{i,j} \), adalah \( (-1)^{i+j} \) kali determinan minor
\( i,j \) dari \( T \).
%</df:Minor>
\end{definition}

\begin{example}
Kofaktor $1,2$ dari matriks dalam persamaan~($*$) adalah negatif dari
determinan pertama dalam~($*\!*\!*$).
\begin{equation*}
  T_{1,2}=
  -1\cdot\begin{vmat}
    t_{2,1}  &t_{2,3}  \\
    t_{3,1}  &t_{3,3}
  \end{vmat}
\end{equation*}
\end{example}

\begin{example}
Untuk
\begin{equation*}
   T=
   \begin{mat}[r]
      1  &2  &3  \\
      4  &5  &6  \\
      7  &8  &9
   \end{mat}
\end{equation*}
berikut adalah kofaktor \( 1,2 \) dan \( 2,2 \).
\begin{equation*}
   T_{1,2}=
   (-1)^{1+2}\cdot\begin{vmat}[r]
                4  &6  \\
                7  &9
             \end{vmat}=6
  \qquad
   T_{2,2}=
   (-1)^{2+2}\cdot\begin{vmat}[r]
                1  &3  \\
                7  &9
             \end{vmat}=-12
\end{equation*}
\end{example}

\begin{theorem}[Ekspansi Laplace untuk Determinan]
\label{th:LaPlaceExp}
\index{determinan!ekspansi Laplace}%
% \index{Laplace expansion!computes determinant}
%<*th:LaPlaceExp>
Jika \( T \) merupakan matriks \( \nbyn{n} \), kita dapat menentukan
determinannya dengan melakukan ekspansi kofaktor pada sebarang
baris~$i$ atau kolom~$j$.
\begin{align*}
   \deter{T}
   &=t_{i,1}\cdot T_{i,1}+t_{i,2}\cdot T_{i,2}+\cdots+t_{i,n}\cdot T_{i,n}  \\
   &=t_{1,j}\cdot T_{1,j}+t_{2,j}\cdot T_{2,j}+\cdots+t_{n,j}\cdot T_{n,j}
\end{align*}
%</th:LaPlaceExp>
\end{theorem}

\begin{proof}
\nearbyexercise{exer:LaplaceProof}.
\end{proof}

\begin{example} \label{ex:ExpLaPlace}
Kita dapat menghitung determinan
\begin{equation*}
   \deter{T}=
   \begin{vmat}[r]
     1  &2  &3  \\
     4  &5  &6  \\
     7  &8  &9
   \end{vmat}
\end{equation*}
dengan melakukan ekspansi sepanjang baris pertama.
\begin{equation*}
   \deter{T}
   =1\cdot(+1)\begin{vmat}[r]
                5  &6  \\
                8  &9
              \end{vmat}
   +2\cdot(-1)\begin{vmat}[r]
                4  &6  \\
                7  &9
              \end{vmat}
   +3\cdot(+1)\begin{vmat}[r]
                4  &5  \\
                7  &8
              \end{vmat}     
  =-3+12-9     
  =0
\end{equation*}
Kita juga dapat melakukan ekspansi sepanjang kolom kedua.
\begin{equation*}
   \deter{T}
   =2\cdot(-1)\begin{vmat}[r]
                4  &6  \\
                7  &9
              \end{vmat}
   +5\cdot(+1)\begin{vmat}[r]
                1  &3  \\
                7  &9
              \end{vmat}
   +8\cdot(-1)\begin{vmat}[r]
                1  &3  \\
                4  &6
              \end{vmat}     
  =12-60+48   
  =0
\end{equation*}
\end{example}

\begin{example}
Baris atau kolom yang memuat banyak nol membuat ekspansi ini lebih mudah.
\begin{equation*}
  \begin{vmat}[r]
    1 &5  &0  \\
    2 &1  &1  \\
    3 &-1 &0
  \end{vmat}
   =
   0\cdot(+1)\begin{vmat}[r]
               2  &1  \\
               3  &-1
             \end{vmat}+
   1\cdot(-1)\begin{vmat}[r]
               1  &5  \\
               3  &-1
             \end{vmat}+
   0\cdot(+1)\begin{vmat}[r]
               1  &5  \\
               2  &1
             \end{vmat}
  =16
\end{equation*}
\end{example}

Sebagai penutup, kita menerapkan ekspansi Laplace untuk menurunkan rumus baru
bagi invers suatu matriks.
Dengan \nearbytheorem{th:LaPlaceExp}, kita dapat menghitung determinan matriks
dengan mengambil kombinasi linear entri-entri suatu baris dan kofaktor yang
bersesuaian.
\begin{equation*}
  t_{i,1}\cdot T_{i,1}+t_{i,2}\cdot T_{i,2}+\dots+t_{i,n}\cdot T_{i,n}
   =\deter{T}  
\end{equation*}
Ingat bahwa matriks dengan dua baris identik memiliki determinan nol.
Karena itu, membobot kofaktor-kofaktor baris~$k$ dengan entri-entri dari
baris~$i$, untuk $k\neq i$, memberikan nol
\begin{equation*}
  t_{i,1}\cdot T_{k,1}+t_{i,2}\cdot T_{k,2}+\dots+t_{i,n}\cdot T_{k,n}=0
\end{equation*}
karena jumlah tersebut merepresentasikan ekspansi sepanjang baris~$k$ dari
matriks yang baris~\( k \)-nya dibuat sama dengan baris~\( i \).
Hal ini dirangkum sebagai berikut.
\begin{equation*}
 \generalmatrix{t}{n}{n}
 \begin{mat}
   T_{1,1}  &T_{2,1}  &\ldots  &T_{n,1}  \\
   T_{1,2}  &T_{2,2}  &\ldots  &T_{n,2}  \\
            &\vdots   &        &         \\
   T_{1,n}  &T_{2,n}  &\ldots  &T_{n,n}
 \end{mat}                                  
 =\begin{mat}
     |T|      &0        &\ldots  &0        \\
     0        &|T|      &\ldots  &0        \\
              &\vdots   &        &         \\
     0        &0        &\ldots  &|T|
   \end{mat} 
\end{equation*}
Perhatikan bahwa urutan subskrip dalam matriks kofaktor berlawanan dengan
urutan subskrip dalam matriks lainnya; sebagai contoh, pada baris pertama
matriks kofaktor, subskripnya adalah $1,1$, lalu $2,1$, dan seterusnya.

\begin{definition}  \label{df:Adjoint}
%<*df:Adjoint>
\definend{Matriks adjoin}\index{matriks!adjoin}\index{matriks adjoin}
(atau \definend{adjugat})\index{matriks!adjugat}\index{matriks adjugat}
dari matriks persegi \( T \)
(juga disebut \definend{adjoin klasik})\index{adjoin klasik}
adalah
\begin{equation*}
  \adj(T)=
    \begin{mat}
      T_{1,1}  &T_{2,1}  &\ldots  &T_{n,1}  \\
      T_{1,2}  &T_{2,2}  &\ldots  &T_{n,2}  \\
               &\vdots   &        &         \\
      T_{1,n}  &T_{2,n}  &\ldots  &T_{n,n}
    \end{mat}
\end{equation*}
di mana entri pada baris~\( i \), kolom~\( j \), yaitu \( T_{j,i} \),
merupakan kofaktor \( j,i \).
%</df:Adjoint>
\end{definition}

\begin{theorem} \label{th:MatTimesAdjEqDiagDets}
%<*th:MatTimesAdjEqDiagDets>
Jika \( T \) merupakan matriks persegi,
$T\cdot \adj(T)=\adj(T)\cdot T=\deter{T}\cdot I$.
Jadi, jika $T$ memiliki invers, yakni jika $\deter{T}\neq 0$, maka
$T^{-1}=(1/\deter{T})\cdot\adj(T)$.\index{matriks!invers}\index{invers!matriks}
%</th:MatTimesAdjEqDiagDets>
\end{theorem}

\begin{proof}
Pembahasan sebelum pernyataan teorema memperjelas hasil tersebut.
\end{proof}

\begin{example}    \label{ex:MatTimesAdjEqDets}
Jika
\begin{equation*}
  T=\begin{mat}[r]
        1  &0  &4  \\
        2  &1  &-1 \\
        1  &0  &1
      \end{mat}
\end{equation*}
maka $\adj(T)$ adalah
\begin{equation*}
    \begin{mat}
      T_{1,1}  &T_{2,1}  &T_{3,1} \\
      T_{1,2}  &T_{2,2}  &T_{3,2} \\
      T_{1,3}  &T_{2,3}  &T_{3,3}   
    \end{mat}
    \!\!=\!\!\begin{mat}
      \begin{vmat}[r]
           1  &-1 \\
           0  &1
         \end{vmat}
      &-\begin{vmat}[r]
             0  &4  \\
             0  &1
           \end{vmat}
      &\begin{vmat}[r]
             0  &4  \\
             1  &-1
         \end{vmat}             \\[2.1ex]
      -\begin{vmat}[r]
           2  &-1 \\
           1  &1
         \end{vmat}
      &\begin{vmat}[r]
             1  &4  \\
             1  &1
           \end{vmat}
      &-\begin{vmat}[r]
             1  &4  \\
             2  &-1
           \end{vmat}            \\[2.1ex]
      \begin{vmat}[r]
           2  &1  \\
           1  &0
         \end{vmat}
      &-\begin{vmat}[r]
             1  &0  \\
             1  &0
           \end{vmat}
      &\begin{vmat}[r]
             1  &0  \\
             2  &1
           \end{vmat}
    \end{mat}
    \!\!=\!                 
    \begin{mat}[r]
       1  &0  &-4  \\
       -3 &-3 &9  \\
       -1 &0  &1
    \end{mat}         
\end{equation*}
dan mengalikannya dengan $T$ menghasilkan matriks diagonal $\deter{T}\cdot I$.
\begin{equation*}
  \begin{mat}[r]
    1  &0  &4  \\
    2  &1  &-1 \\
    1  &0  &1
  \end{mat}
  \begin{mat}[r]
    1  &0  &-4  \\
    -3 &-3 &9  \\
    -1 &0  &1
   \end{mat}         
  =\begin{mat}[r]
     -3  &0  &0  \\
      0  &-3 &0  \\
      0  &0  &-3
   \end{mat}
\end{equation*}
% \end{example}

% \begin{corollary} \label{cor:InvFromAdj}
% \index{matrix!inverse}\index{inverse!matrix}
% %<*co:InvFromAdj>
% If \( \deter{T}\neq 0 \) then
% $T^{-1}=(1/\deter{T})\cdot\adj(T)$.
% %</co:InvFromAdj>
% \end{corollary}

% \begin{example}
\noindent Invers dari $T$
% The inverse of the matrix from \nearbyexample{ex:MatTimesAdjEqDets}
adalah $(1/-3)\cdot\adj(T)$.
\begin{equation*}
  T^{-1}
  =\begin{mat}[r]  % braces make the - a negative sign?
     1/\hbox{$-3$}  &0/\hbox{$-3$}  &-4/\hbox{$-3$}  \\
    -3/\hbox{$-3$}  &-3/\hbox{$-3$} &9/\hbox{$-3$}   \\
    -1/\hbox{$-3$}  &0/\hbox{$-3$}  &1/\hbox{$-3$}
   \end{mat}
  =\begin{mat}[r]
     -1/3  &0  &4/3  \\
      1    &1  &-3   \\
      1/3  &0  &-1/3
   \end{mat}
\end{equation*}
\end{example}

Rumus-rumus dari subbagian ini sering digunakan untuk perhitungan dengan tangan
dan kadang-kadang berguna bagi jenis matriks khusus.
Namun, bagi matriks umum, rumus-rumus ini bukan pilihan terbaik karena
memerlukan lebih banyak aritmetika daripada, misalnya, metode Gauss--Jordan.




\begin{exercises}
  \recommended \item 
    Tentukan kofaktor berikut.
    \begin{equation*}
      T=\begin{mat}[r]
          1  &0  &2  \\
         -1  &1  &3  \\
          0  &2  &-1
        \end{mat}
    \end{equation*}
    \begin{exparts*}
       \partsitem \( T_{2,3} \)
       \partsitem \( T_{3,2} \)
       \partsitem \( T_{1,3} \)
    \end{exparts*}
    \begin{answer}
      \begin{exparts*}
       \partsitem \( (-1)^{2+3}\begin{vmat}[r]
                            1  &0  \\
                            0  &2
                          \end{vmat}=-2  \)
       \partsitem \( (-1)^{3+2}\begin{vmat}[r]
                            1  &2  \\
                           -1  &3
                          \end{vmat}=-5  \)
       \partsitem \( (-1)^{4}\begin{vmat}[r]
                           -1  &1  \\
                            0  &2
                          \end{vmat}=-2  \)
      \end{exparts*}  
    \end{answer}
  \recommended \item 
    Tentukan matriks adjoin dari matriks berikut.
    \begin{equation*}
      T=\begin{mat}[r]
          1  &0  &2  \\
         -1  &1  &3  \\
          0  &2  &-1
        \end{mat}
    \end{equation*}
    \begin{answer}
      % sage: M = matrix(QQ, [[1,0,2], [-1,1,3], [0,2,-1]])
      % sage: M.adjoint()
      % [-7  4 -2]
      % [-1 -1 -5]
      % [-2 -2  1] 
      {\renewcommand{\arraystretch}{1.2}
      \begin{equation*} 
        \adj(T)=
        \begin{mat}
          \begin{vmatrix}
            1 &3 \\
            2 &-1
          \end{vmatrix}
          &-\begin{vmatrix}
            0 &2 \\ 
            2 &-1
          \end{vmatrix}
          &\begin{vmatrix}
             0 &2 \\
              1 &3
           \end{vmatrix}                  \\
          -\begin{vmatrix}
             -1 &3 \\
             0 &-1
           \end{vmatrix}
          &\begin{vmatrix}
            1 &2 \\
            0 &-1
           \end{vmatrix}
          &-\begin{vmatrix}
            1 &2 \\
            -1 &3
           \end{vmatrix}                  \\
          \begin{vmatrix}
            -1 &1 \\ 
            0 &2
          \end{vmatrix}
          &-\begin{vmatrix}
            1  &0 \\
            0 &2
           \end{vmatrix}
          &\begin{vmatrix}
            1 &0 \\
            -1 &1
           \end{vmatrix}                  
        \end{mat}
        =
        \begin{mat}
          -7 &4 &-2 \\
          -1  &-1 &-5 \\
          -2 &-2 &1 
        \end{mat}
      \end{equation*}}
    \end{answer}
  \item Determinan ini bernilai $0$.
   Hitunglah nilainya dengan melakukan ekspansi pada baris pertama.
    \begin{equation*}
      \begin{vmat}[r]
         1  &2  &3  \\
         4  &5  &6  \\
         7  &8  &9
      \end{vmat}
    \end{equation*}
    \begin{answer}
      \begin{equation*}
        \begin{vmat}[r]
          1  &2  &3  \\
          4  &5  &6  \\
          7  &8  &9
        \end{vmat}=
        1\cdot(+1)\cdot
          \begin{vmatrix}
            5 &6 \\
            8 &9
          \end{vmatrix}
        +2\cdot(-1)\cdot
          \begin{vmatrix}
            4 &6 \\
            7 &9
          \end{vmatrix}
        +3\cdot(+1)\cdot
          \begin{vmatrix}
            4 &5 \\
            7 &8
          \end{vmatrix}
      \end{equation*}
      Rumus bagi matriks $\nbyn{2}$ memberikan
      $1\cdot (-3)-2\cdot (-6)+3\cdot(-3)=0$.
    \end{answer}
  \recommended \item 
    Tentukan determinan berikut dengan melakukan ekspansi
    \begin{equation*}
      \begin{vmat}[r]
         3  &0  &1  \\
         1  &2  &2  \\
        -1  &3  &0
      \end{vmat}
    \end{equation*}
    \begin{exparts*}
      \partsitem pada baris pertama
      \partsitem pada baris kedua
      \partsitem pada kolom ketiga.
    \end{exparts*}
    \begin{answer}
      \begin{exparts}
         \partsitem \( 3\cdot(+1)\begin{vmat}[r]
                          2  &2  \\
                          3  &0
                       \end{vmat}
                 +0\cdot (-1)\begin{vmat}[r]
                          1  &2  \\
                         -1  &0
                       \end{vmat}
                 +1\cdot (+1)\begin{vmat}[r]
                          1  &2  \\
                         -1  &3
                       \end{vmat} =-13 \)
         \partsitem \( 1\cdot (-1)\begin{vmat}[r]
                          0  &1  \\
                          3  &0
                       \end{vmat}
                 +2\cdot (+1)\begin{vmat}[r]
                          3  &1  \\
                         -1  &0
                       \end{vmat}
                 +2\cdot (-1)\begin{vmat}[r]
                          3  &0  \\
                         -1  &3
                       \end{vmat} =-13 \)
         \partsitem \( 1\cdot (+1)\begin{vmat}[r]
                          1  &2  \\
                         -1  &3
                       \end{vmat}
                 +2\cdot (-1)\begin{vmat}[r]
                          3  &0  \\
                         -1  &3
                       \end{vmat}
                 +0\cdot (+1)\begin{vmat}[r]
                          3  &0  \\
                          1  &2
                       \end{vmat} =-13 \)
       \end{exparts}  
     \end{answer}
  \item 
    Tentukan matriks adjoin dari matriks dalam \nearbyexample{ex:ExpLaPlace}.
    \begin{answer}
      Berikut adalah $\adj(T)$.
      \begin{align*}
         \begin{mat}
           T_{1,1}  &T_{2,1}  &T_{3,1} \\
           T_{1,2}  &T_{2,2}  &T_{3,2} \\
           T_{1,3}  &T_{2,3}  &T_{3,3}   
         \end{mat}
         &=\begin{mat}
            \begin{vmat}
               5  &6  \\ 8  &9
            \end{vmat}
            &-\begin{vmat}
               2  &3  \\  8  &9
            \end{vmat}
            &+\begin{vmat}
               2  &3  \\  5  &6
            \end{vmat}        \\[2.5ex]
            -\begin{vmat}
               4  &6  \\  7  &9
            \end{vmat}
            &+\begin{vmat}
               1  &3  \\  7  &9
            \end{vmat}
            &-\begin{vmat}
               1  &3  \\  4  &6
            \end{vmat}        \\[2.5ex]
            +\begin{vmat}
               4  &5  \\  7  &8
            \end{vmat}
            &-\begin{vmat}
               1  &2  \\  7  &8
            \end{vmat}
            &+\begin{vmat}
               1  &2  \\  4  &5
            \end{vmat}
         \end{mat}                                   \\
         &=\begin{mat}[r]
          -3  &6   &-3 \\
           6  &-12 &6  \\
          -3  &6   &-3
         \end{mat}
      \end{align*}
    \end{answer}
  \recommended \item 
    Tentukan matriks adjoin dari masing-masing matriks berikut.
    \begin{exparts*}
      \partsitem \( \begin{mat}[r]
                 2   &1  &4  \\
                -1   &0  &2  \\
                 1   &0  &1
               \end{mat}  \)
      \partsitem \( \begin{mat}[r]
                 3  &-1  \\
                 2  &4
               \end{mat}  \)
      \partsitem \( \begin{mat}[r]
                 1   &1  \\
                 5   &0
               \end{mat}  \)
      \partsitem \( \begin{mat}[r]
                 1   &4  &3  \\
                -1   &0  &3  \\
                 1   &8  &9
               \end{mat}  \)
    \end{exparts*}
    \begin{answer}
      \begin{exparts}
        \partsitem Berikut adalah matriks adjoin tersebut.
          \begin{align*}
          \begin{mat}
            T_{1,1}  &T_{2,1}  &T_{3,1} \\ 
            T_{1,2}  &T_{2,2}  &T_{3,2} \\ 
            T_{1,3}  &T_{2,3}  &T_{3,3} 
          \end{mat}
          &=\begin{mat}
            \begin{vmat}[r]
              0  &2  \\
              0  &1 
            \end{vmat}
            &-\begin{vmat}[r]
               1  &4  \\
               0  &1
            \end{vmat}
            &\begin{vmat}[r]
               1  &4  \\
               0  &2 
            \end{vmat}        \\[2.5ex]
            -\begin{vmat}[r]
              -1  &2  \\
               1  &1
            \end{vmat}
            &\begin{vmat}[r]
               2  &4  \\
               1  &1
            \end{vmat}
            &-\begin{vmat}[r]
               2  &4  \\ 
              -1  &2
            \end{vmat}         \\[2.5ex]
            \begin{vmat}[r]
              -1  &0  \\
               1  &0
            \end{vmat} 
            &-\begin{vmat}[r]
               2  &1  \\
               1  &0
            \end{vmat}
            &\begin{vmat}[r]
               2  &1  \\
              -1  &0
            \end{vmat}
          \end{mat}                         \\
          &=\begin{mat}[r]
             0  &-1  &2  \\
             3  &-2  &-8 \\
             0  &1   &1            
          \end{mat}
          \end{align*}
        \partsitem Minor-minornya berukuran $\nbyn{1}$.
          $\begin{mat}
             T_{1,1}  &T_{2,1} \\
             T_{1,2}  &T_{2,2}  
          \end{mat}
          =
          \begin{mat}
            \begin{vmat}
              4              
            \end{vmat}
            &-\begin{vmat}
              -1
            \end{vmat}        \\[1.5ex]
            -\begin{vmat}
               2              
            \end{vmat}
            &\begin{vmat}
               3
            \end{vmat}
          \end{mat}
          =
          \begin{mat}[r]
            4  &1  \\
           -2  &3
          \end{mat}$
        \partsitem 
          $\begin{mat}[r]
               0  &-1 \\
              -5  &1
          \end{mat}$
        \partsitem Minor-minornya berukuran $\nbyn{2}$.
          \begin{align*}
          \begin{mat}
            T_{1,1}  &T_{2,1}  &T_{3,1} \\ 
            T_{1,2}  &T_{2,2}  &T_{3,2} \\ 
            T_{1,3}  &T_{2,3}  &T_{3,3} 
          \end{mat}
          &=\begin{mat}
            \begin{vmat}[r]
              0  &3  \\
              8  &9
            \end{vmat}
            &-\begin{vmat}[r]
              4  &3  \\
              8  &9
            \end{vmat}
            &\begin{vmat}[r]
              4  &3  \\
              0  &3 
            \end{vmat}       \\[2.5ex]
            -\begin{vmat}[r]
              -1  &3  \\
               1  &9
            \end{vmat}
            &\begin{vmat}[r]
               1  &3  \\
               1  &9              
            \end{vmat}
            &-\begin{vmat}[r]
               1  &3 \\
              -1  &3              
            \end{vmat}        \\[2.5ex]
            \begin{vmat}[r]
               -1  &0  \\
                1  &8              
            \end{vmat}
            &-\begin{vmat}[r]
                1  &4 \\
                1  &8              
            \end{vmat}
            &\begin{vmat}[r]
                1  &4  \\
               -1  &0              
            \end{vmat}
          \end{mat}                       \\
          &=
          \begin{mat}[r]
            -24  &-12  &12  \\
             12  &6    &-6   \\
             -8  &-4   &4
          \end{mat}
        \end{align*}
      \end{exparts}
    \end{answer}
\recommended \item
    Tentukan invers setiap matriks pada soal sebelumnya dengan menggunakan
    \nearbytheorem{th:MatTimesAdjEqDiagDets}.
    \begin{answer}
      \begin{exparts}
        \partsitem $(1/3)\cdot 
           \begin{mat}[r]
             0  &-1  &2  \\
             3  &-2  &-8 \\
             0  &1   &1            
          \end{mat}          
          =
          \begin{mat}[r]
            0  &-1/3  &2/3  \\
            1  &-2/3  &-8/3 \\ 
            0  &1/3   &1/3
          \end{mat}$
        \partsitem $(1/14)\cdot 
          \begin{mat}[r]
            4  &1  \\
           -2  &3
          \end{mat}
          =
          \begin{mat}[r]
            2/7  &1/14  \\
           -1/7  &3/14
          \end{mat}$
        \partsitem $(1/-5)\cdot 
          \begin{mat}[r]
               0  &-1 \\
              -5  &1
          \end{mat}
          =
          \begin{mat}[r]
            0  &1/5 \\
            1  &-1/5
          \end{mat}$
        \partsitem Matriks tersebut memiliki determinan nol, sehingga tidak
           memiliki invers.
      \end{exparts}
    \end{answer}
  \item 
    Tentukan matriks adjoin dari matriks berikut.
    \begin{equation*}
      \begin{mat}[r]
        2  &1  &0  &0  \\
        1  &2  &1  &0  \\
        0  &1  &2  &1  \\
        0  &0  &1  &2
      \end{mat}
    \end{equation*}
    \begin{answer}
      $\begin{mat}
        T_{1,1}  &T_{2,1}  &T_{3,1}  &T_{4,1}  \\ 
        T_{1,2}  &T_{2,2}  &T_{3,2}  &T_{4,2}  \\ 
        T_{1,3}  &T_{2,3}  &T_{3,3}  &T_{4,3}  \\ 
        T_{1,4}  &T_{2,4}  &T_{3,4}  &T_{4,4}  
      \end{mat}
      =
      \begin{mat}[r]
        4  &-3  &2  &-1  \\
       -3  &6   &-4 &2   \\
        2  &-4  &6  &-3  \\
        -1 &2  &-3  &4
      \end{mat}$
    \end{answer}
  \recommended \item
    Lakukan ekspansi sepanjang baris pertama untuk menurunkan rumus determinan
    matriks \( \nbyn{2} \).
    \begin{answer}
       Determinan
       \begin{equation*}
         \begin{vmat}
           a  &b  \\
           c  &d
         \end{vmat}
       \end{equation*} 
       yang diekspansi sepanjang baris pertama memberikan
       \( a\cdot (+1)\deter{d}+b\cdot (-1)\deter{c}=ad-bc \)
       (perhatikan kedua minor $\nbyn{1}$).
     \end{answer}
  \recommended \item
    Lakukan ekspansi sepanjang baris pertama untuk menurunkan rumus determinan
    matriks \( \nbyn{3} \).
    \begin{answer}
      Determinan dari
      \begin{equation*}
        \begin{mat}
          a  &b  &c  \\
          d  &e  &f  \\
          g  &h  &i
        \end{mat}
      \end{equation*}
      adalah sebagai berikut.
      \begin{equation*}
        a\cdot \begin{vmat}
            e  &f  \\
            h  &i
         \end{vmat}
        -b\cdot \begin{vmat}
            d  &f  \\
            g  &i
         \end{vmat}
        +c\cdot \begin{vmat}
            d  &e  \\
            g  &h
         \end{vmat}
        =a(ei-fh)-b(di-fg)+c(dh-eg)
      \end{equation*}  
    \end{answer}
  \recommended \item 
   \begin{exparts}
     \partsitem Berikan rumus bagi matriks adjoin dari matriks \( \nbyn{2} \).
     \partsitem Gunakan rumus tersebut untuk menurunkan rumus invers.
   \end{exparts} 
    \begin{answer}
      \begin{exparts}
        \partsitem
          $\begin{mat}
             T_{1,1}  &T_{2,1}  \\
             T_{1,2}  &T_{2,2}
            \end{mat}
            =\begin{mat}
            \begin{vmat}
              t_{2,2}
            \end{vmat}
           &-\begin{vmat}
               t_{1,2}
             \end{vmat}       \\
           -\begin{vmat}
               t_{2,1}
            \end{vmat}
           &\begin{vmat}
              t_{1,1} 
            \end{vmat}
          \end{mat}
          =\begin{mat}
             t_{2,2}  &-t_{1,2}  \\
            -t_{2,1} &t_{1,1}
           \end{mat}$
        \partsitem
          $\bigl(1/(t_{1,1}t_{2,2}-t_{1,2}t_{2,1})\bigr)\cdot
          \begin{mat}
             t_{2,2}  &-t_{1,2}  \\
            -t_{2,1} &t_{1,1}
           \end{mat}$
      \end{exparts}
    \end{answer}
  \recommended \item
    Dapatkah kita menghitung determinan dengan melakukan ekspansi sepanjang
    diagonal?
    \begin{answer}
      Tidak.
      Berikut sebuah determinan yang nilainya
      \begin{equation*}
        \begin{vmat}[r]
          1  &0  &0  \\
          0  &1  &0  \\
          0  &0  &1
        \end{vmat}=1
      \end{equation*}
      tidak sama dengan hasil ekspansi sepanjang diagonal.
      \begin{equation*}
        1\cdot (+1)\begin{vmat}[r]
               1  &0  \\
               0  &1
             \end{vmat}
       +1\cdot (+1)\begin{vmat}[r]
               1  &0  \\
               0  &1
             \end{vmat}
       +1\cdot (+1)\begin{vmat}[r]
               1  &0  \\
               0  &1
             \end{vmat}=3
      \end{equation*}  
    \end{answer}
  \item 
    Berikan rumus bagi matriks adjoin dari suatu matriks diagonal.
    \begin{answer}
      Tinjau matriks diagonal berikut.
      \begin{equation*}
        D=
        \begin{mat}
          d_1  &0   &0   &\ldots    \\
          0    &d_2 &0   &          \\
          0    &0   &d_3            \\
               &    &    &\ddots    \\
               &    &    &      &d_n   
        \end{mat}
      \end{equation*}
      Jika $i\neq j$, minor $i,j$ merupakan matriks $\nbyn{(n-1)}$ yang hanya
      memiliki $n-2$ entri tak nol, karena kita telah menghapus $d_i$ dan
      $d_j$.
      Jadi, sekurang-kurangnya satu baris atau kolom minor tersebut seluruhnya
      nol, sehingga kofaktor $D_{i,j}$ bernilai nol.
      Jika $i=j$, minornya adalah matriks diagonal dengan entri $d_1$, \ldots,
      $d_{i-1}$, $d_{i+1}$, \ldots, $d_n$.
      Determinan minor tersebut adalah hasil kali semua entri itu, sedangkan
      faktor tanda kofaktornya ialah $(-1)^{i+j}=(-1)^{2i}=1$.
      \begin{equation*}
        \adj(D)
        =
        \begin{mat}
          d_2\cdots d_n    &0                 &      &0    \\
          0                &d_1d_3\cdots d_n  &      &0    \\
                           &                  &\ddots      \\
                           &                  &       &d_1\cdots d_{n-1} 
        \end{mat}
      \end{equation*}

      Sebagai tambahan, \nearbytheorem{th:MatTimesAdjEqDiagDets} memberikan
      cara yang lebih ringkas untuk menurunkan kesimpulan ini.
    \end{answer}
  \recommended \item
    Buktikan bahwa transpos dari matriks adjoin sama dengan matriks adjoin dari
    transposnya.
    \begin{answer}
      Perhatikan saja bahwa jika $S=\trans{T}$, maka kofaktor $S_{j,i}$ sama
      dengan kofaktor $T_{i,j}$ karena $(-1)^{j+i}=(-1)^{i+j}$ dan karena
      kedua minornya saling merupakan transpos (sedangkan determinan transpos
      sama dengan determinan matriks semula).
    \end{answer}
  \item 
    Buktikan atau sangkal: \( \adj(\adj(T))=T \).
    \begin{answer}
      Pernyataan ini salah; berikut sebuah contoh.
      \begin{equation*}
         T=\begin{mat}[r]
           1  &2   &3  \\
           4  &5   &6  \\
           7  &8   &9
         \end{mat}
         \qquad
         \adj(T)=\begin{mat}[r]
          -3  &6   &-3 \\
           6  &-12 &6  \\
          -3  &6   &-3
         \end{mat}
         \qquad
         \adj(\adj(T))=\begin{mat}[r]
           0  &0   &0  \\
           0  &0   &0  \\
           0  &0   &0
         \end{mat}
      \end{equation*}
    \end{answer}
\item
    Sebuah matriks persegi disebut \definend{segitiga atas}
    \index{matriks!segitiga} jika setiap entri \( i,j \) pada bagian di bawah
    diagonal bernilai nol, yakni ketika \( i>j \).
    \begin{exparts}
      \partsitem
        Apakah matriks adjoin dari matriks segitiga atas harus berbentuk
        segitiga atas? Segitiga bawah?
      \partsitem Buktikan bahwa invers suatu matriks segitiga atas juga
        berbentuk segitiga atas, jika invers tersebut ada.
    \end{exparts}
    \begin{answer}
      \begin{exparts}
        \partsitem Sebuah contoh
          \begin{equation*}
             M=
            \begin{mat}[r]
              1  &2  &3  \\
              0  &4  &5  \\
              0  &0  &6
            \end{mat}
          \end{equation*}
          mengisyaratkan jawaban yang benar.
          \begin{align*}
            \adj(M)=
            \begin{mat}
             M_{1,1}  &M_{2,1}  &M_{3,1}  \\
             M_{1,2}  &M_{2,2}  &M_{3,2}  \\
             M_{1,3}  &M_{2,3}  &M_{3,3}  
            \end{mat}
            &=\begin{mat}
              \begin{vmat}
                4  &5 \\ 0 &6
              \end{vmat}
              &-\begin{vmat}
                2  &3  \\  0  &6
              \end{vmat}
              &\begin{vmat}
                2  &3  \\  4  &5 
              \end{vmat}          \\[1.5ex]
              -\begin{vmat}
                0  &5  \\  0  &6
              \end{vmat}
              &\begin{vmat}
                1  &3  \\  0  &6
              \end{vmat}
              &-\begin{vmat}
                1  &3  \\  0  &5  
              \end{vmat}           \\[1.5ex]
              \begin{vmat}
                0  &4  \\  0  &0
              \end{vmat}
              &-\begin{vmat}
                1  &2  \\  0  &0
              \end{vmat}
              &\begin{vmat}
                1  &2  \\  0  &4
              \end{vmat}
            \end{mat}                                \\
            &=
            \begin{mat}[r]
              24  &-12 &-2 \\
               0  &6   &-5  \\
               0  &0   &4
            \end{mat}
          \end{align*}
          Hasilnya memang berbentuk segitiga atas.

          Verifikasi ini terperinci, tetapi tidak sulit.
          Entri-entri di bawah diagonal matriks adjoin adalah $M_{a,b}$ dengan
          $a<b$.
          Kita perlu memverifikasi bahwa kofaktor $M_{a,b}$ bernilai nol jika
          $a<b$.
          Dengan $a<b$, hapus baris~$a$ dan kolom~$b$ dari $M$.
          \begin{equation*}
            \begin{mat}
              m_{1,1} &\ldots   &m_{1,b}  &\ldots       &        \\
              m_{2,1} &\ldots   &m_{2,b}  &       &        \\
              \vdots  &         &\vdots   &       &        \\
              m_{a,1} &\ldots   &m_{a,b}  &\ldots &m_{a,n} \\
              \vdots        &         &\vdots   &       &        \\
                      &         &m_{n,b}  &       &                   
            \end{mat}
          \end{equation*}
          setelah dihapus, menyisakan minor berbentuk segitiga atas.
          Entri~$i,j$ dari minor tersebut dapat berupa entri~$i,j$ dari $M$
          (jika $i<a$ dan~$j<b$; dalam kasus ini, $i>j$ menyiratkan bahwa
          entri tersebut nol), atau entri~$i,j+1$ dari $M$ (jika $i<a$ dan
          $j\geq b$; posisi di bawah diagonal tidak mungkin terjadi karena
          $i<a<b\leq j$), atau entri~$i+1,j$ dari $M$ (jika $i\geq a$ dan
          $j<b$; dalam kasus ini, $i>j$ menyiratkan $i+1>j$), atau entri
          $i+1,j+1$ dari $M$ (jika $i\geq a$ dan~$j\geq b$; di sini $i>j$
          menyiratkan $i+1>j+1$). Dengan demikian, setiap entri di bawah
          diagonal minor bernilai nol.
          Determinan minor tersebut merupakan hasil kali entri-entri diagonalnya.
          Perhatikan bahwa entri $a+1,a$ dari $M$ menjadi entri
          $a,a$ dari minor (entri itu tidak terhapus karena $a<b$ secara ketat).
          Entri tersebut nol karena $M$ berbentuk segitiga atas dan $a+1>a$.
          Jadi, kofaktornya nol dan matriks adjoin berbentuk segitiga atas.
          (Kasus segitiga bawah serupa.)
        \partsitem Hasil ini langsung mengikuti butir sebelumnya dan
          \nearbytheorem{th:MatTimesAdjEqDiagDets}.
      \end{exparts}
    \end{answer}
  \item 
    \label{exer:LaplaceProof}
    \textit{Soal ini memerlukan materi dari subbagian opsional Keberadaan
       Determinan.}
    Buktikan \nearbytheorem{th:LaPlaceExp} dengan menggunakan ekspansi
    permutasi.
    \begin{answer}
      Kita akan menunjukkan bahwa setiap determinan dapat diekspansi sepanjang
      baris~\( i \).
      Argumen bagi kolom~$j$ serupa.

      Setiap suku dalam ekspansi permutasi memuat tepat satu entri dari setiap
      baris.
      Seperti dalam \nearbyexample{ex:ExpLaPlace}, keluarkan setiap entri
      baris~$i$ sebagai faktor untuk memperoleh
      \( \deter{T}
           =t_{i,1}\cdot\hat{T}_{i,1}+\dots+t_{i,n}\cdot\hat{T}_{i,n} \),
      dengan setiap \( \hat{T}_{i,j} \) merupakan jumlah suku-suku yang tidak
      memuat satu pun entri dari baris \( i \).
      Kita akan menunjukkan bahwa \( \hat{T}_{i,j} \) adalah kofaktor
      \( i,j \).

      Pertama-tama, tinjau kasus \( i,j=n,n \):
      \begin{equation*}
        t_{n,n}\cdot\hat{T}_{n,n}
        =t_{n,n}\cdot
            \sum_{\phi}t_{1,\phi(1)}t_{2,\phi(2)}\dots\,t_{n-1,\phi(n-1)}
                      \sgn(\phi)
      \end{equation*}
      dengan penjumlahan dilakukan atas semua permutasi berukuran \( n \),
      \( \phi \), yang
      memenuhi \( \phi(n)=n \).
      Untuk menunjukkan bahwa \( \hat{T}_{n,n} \) sama dengan kofaktor
      \( T_{n,n} \), kita hanya perlu menunjukkan bahwa jika \( \phi \)
      merupakan permutasi berukuran \( n \) dengan \( \phi(n)=n \), sedangkan
      \( \sigma \) merupakan permutasi berukuran \( n-1 \) dengan
      \( \sigma(1)=\phi(1) \), \ldots,
      \( \sigma(n-1)=\phi(n-1) \), maka
      \( \sgn(\sigma)=\sgn(\phi) \).
      Hal ini benar karena $\phi$ dan $\sigma$ memiliki banyak inversi yang sama.

      Kembali ke kasus umum \( i,j \).
      Tukarkan baris-baris yang bersebelahan hingga baris ke-\( i \) menjadi
      baris terakhir, lalu tukarkan kolom-kolom yang bersebelahan hingga kolom
      ke-\( j \) menjadi kolom terakhir.
      Determinan minor \( i,j \) tidak terpengaruh oleh pertukaran-pertukaran
      ini karena baris ke-\( i \) dan kolom ke-\( j \) memang dihilangkan dari
      minor tersebut.
      Sebaliknya, tanda \( \deter{T} \), dan karenanya tanda
      \( \hat{T}_{i,j} \), berubah sebanyak \( n-i+n-j \) kali.
      Nyatakan minor \( i,j \) dengan \( M_{i,j} \). Jadi,
      \( \hat{T}_{i,j}=(-1)^{n-i+n-j}\deter{M_{i,j}}
                =(-1)^{i+j}\deter{M_{i,j}}=T_{i,j} \),
      sehingga \( \hat{T}_{i,j} \) tepat sama dengan kofaktor \( T_{i,j} \).
    \end{answer}
  \item 
    Dengan menggunakan ekspansi Laplace dan induksi pada ukuran matriks,
    buktikan bahwa determinan suatu matriks sama dengan determinan transposnya.
    \begin{answer}
      Pernyataan ini jelas bagi kasus dasar \( \nbyn{1} \).

      Untuk kasus induksi, andaikan determinan suatu matriks sama dengan
      determinan transposnya bagi semua matriks \( \nbyn{1} \), \ldots,
      \( \nbyn{(n-1)} \).
      Ekspansi sepanjang baris~\( i \) memberikan
      \( \deter{T}=t_{i,1}T_{i,1}+\dots\,+t_{i,n}T_{i,n} \)
      sedangkan ekspansi sepanjang kolom~\( i \) memberikan
      \( \deter{\trans{T}}=t_{i,1}(\trans{T})_{1,i}
           +\dots+t_{i,n}(\trans{T})_{n,i} \).
      Karena \( (-1)^{i+j}=(-1)^{j+i} \), tanda-tandanya sama dalam kedua
      penjumlahan.
      Minor \( j,i \) dari \( \trans{T} \) merupakan transpos dari minor
      \( i,j \) dari \( T \); berdasarkan hipotesis induksi, kedua minor itu
      memiliki determinan yang sama. Dengan demikian,
      \( (\trans{T})_{j,i}=T_{i,j} \), sehingga kedua ekspansi di atas sama.
    \end{answer}
  \puzzle \item 
    Tunjukkan bahwa
    \begin{equation*}
      F_n=
      \begin{vmat}
        1  &-1  &1  &-1  &1  &-1  &\ldots  \\
        1  &1   &0  &1   &0  &1   &\ldots  \\
        0  &1   &1  &0   &1  &0   &\ldots  \\
        0  &0   &1  &1   &0  &1   &\ldots  \\
        .  &.   &.  &.   &.  &.   &\ldots
      \end{vmat}
    \end{equation*}
    dengan \( F_n \) merupakan suku ke-\( n \) dari barisan Fibonacci
    \( 1,1,2,3,5,\dots,x,y,x+y,\ldots\, \), sedangkan determinannya berordo
    \( n-1 \).
    \cite{Monthly49p409}
    \begin{answer}
      \answerasgiven %
      Nyatakan determinan di atas dengan \( D_n \). Jelas bahwa
      \( D_2=1 \) dan \( D_3=2 \).
      Tinggal menunjukkan bahwa \( D_n=D_{n-1}+D_{n-2},\; n\geq 4 \).
      Dalam \( D_n \), kurangkan kolom ke-\( (n-3) \) dari kolom ke-\( (n-1) \),
      kolom ke-\( (n-4) \) dari kolom ke-\( (n-2) \), \ldots, dan kolom
      pertama dari kolom ketiga, sehingga diperoleh
      \begin{equation*}
        D_n=
        \begin{vmat}
          1  &-1  &0  &0   &0  &0   &\ldots  \\
          1  &1   &-1 &0   &0  &0   &\ldots  \\
          0  &1   &1  &-1  &0  &0   &\ldots  \\
          0  &0   &1  &1   &-1 &0   &\ldots  \\
          .  &.   &.  &.   &.  &.   &\ldots
        \end{vmat}.
      \end{equation*}
      Ekspansi determinan ini sepanjang baris pertama menghasilkan relasi yang
      diinginkan.
    \end{answer}
%  \recommended \item 
%    Prove that a square matrix is singular if and only if its adjoint is
%    singular.
%  \item 
%    Prove that if \( T \) is an \( \nbyn{n} \) matrix then
%    \( \det(\adj(T))=(\det(T))^{n-1} \).
%    \begin{answer}
%      By \nearbytheorem{th:MatTimesAdjEqDiagDets}, 
%      $\deter{\adj(T)\cdot T}=\deter{T}^n$ because 
%      $\deter{T}\cdot I$ is a diagonal matrix.
%      Then $\deter{T}^n=\deter{\adj(T)\cdot T}
%                         =\deter{\adj(T)}\cdot \deter{T}$.
%    \end{answer}
\end{exercises}
\index{ekspansi determinan Laplace|)}
